% !TeX program = xelatex
% Numerical tables are generated by refresh_tables.py.
\documentclass[10pt,a4paper]{article}
\usepackage[top=17mm,bottom=18mm,left=18mm,right=18mm]{geometry}
\usepackage{fontspec}
\setmainfont{TeX Gyre Termes}
\setsansfont{TeX Gyre Heros}
\setmonofont{TeX Gyre Cursor}[Scale=0.82]
\usepackage{amsmath,amssymb,booktabs,array,tabularx,longtable}
\usepackage[table]{xcolor}
\usepackage{graphicx,fancyhdr,lastpage,enumitem}
\usepackage[unicode,colorlinks=true,linkcolor=teal!65!black,urlcolor=teal!65!black]{hyperref}
\definecolor{ink}{HTML}{16313F}
\definecolor{accent}{HTML}{176B70}
\definecolor{pale}{HTML}{EDF5F5}
\definecolor{muted}{HTML}{596874}
\definecolor{waiting}{HTML}{805D29}
\hypersetup{pdftitle={ISLR - Báo cáo ablation study F0-F4},pdfauthor={ISLR Project},pdfsubject={Feature ablation, validation và protocol tái lập}}
\pagestyle{fancy}
\fancyhf{}
\fancyhead[L]{\small\sffamily\color{ink} ISLR / ABLATION STUDY}
\fancyhead[R]{\small\color{muted} Cập nhật \ReportDate}
\fancyfoot[L]{\small\color{muted} Feature campaign / validation}
\fancyfoot[R]{\small Trang \thepage/\pageref{LastPage}}
\setlength{\headheight}{14pt}
\setlength{\parindent}{0pt}
\setlength{\parskip}{5pt}
\raggedbottom
\makeatletter
\renewcommand\section{\@startsection{section}{1}{\z@}{-2.1ex plus -.5ex minus -.2ex}{1.0ex plus .2ex}{\normalfont\large\bfseries\color{ink}}}
\makeatother
\renewcommand{\arraystretch}{1.22}
\setlength{\tabcolsep}{4pt}
\setlist[itemize]{leftmargin=14pt,itemsep=2pt,topsep=3pt}
\setlist[enumerate]{leftmargin=17pt,itemsep=3pt,topsep=3pt}
\newcolumntype{L}[1]{>{\raggedright\arraybackslash}p{#1}}
\newcolumntype{Y}{>{\raggedright\arraybackslash}X}
\newcommand{\pending}{\textcolor{waiting}{\textsf{PENDING}}}
\newcommand{\code}[1]{\texttt{#1}}
\newcommand{\pageheading}[2]{\vspace*{1mm}{\large\sffamily\bfseries\color{accent}#1}\par{\LARGE\sffamily\bfseries\color{ink}#2}\par\vspace{2mm}}
\newcommand{\callout}[1]{\par\vspace{2mm}\noindent\colorbox{pale}{\parbox{\dimexpr\linewidth-2\fboxsep\relax}{#1}}\par\vspace{2mm}}
\input{data.tex}
\input{tables/summary.tex}
\input{tables/comparisons.tex}
\input{tables/seeds.tex}
\input{tables/paired.tex}
\input{tables/signers.tex}
\input{tables/glosses.tex}
\input{tables/cohort.tex}
\begin{document}

\pageheading{BÁO CÁO THỰC NGHIỆM / 01}{Ablation study cho ISLR}
\textbf{Project:} SignLanguageReconition\_ML. \textbf{Phạm vi:} sáu cấu hình F0, F0-control, F1, F2, F3, F4; mỗi cấu hình train độc lập với seeds 42/43/44. Báo cáo ghi lý do chạy, thay đổi cấu hình, đầu ra và bằng chứng validation.

\callout{\textbf{Kết quả hiện có.} F0 đạt Macro-F1 \BaselineFone\%; F0-control đạt \ControlFone\%, chênh lệch \FoneDelta{} điểm phần trăm. Đã nhận đủ \CompletedConfigs/6 cấu hình, tương ứng \CompletedModels/18 model. \PendingSummary{} Chưa chốt cấu hình thắng toàn campaign.}

\section*{Bảng 1. Kết quả validation của toàn campaign}
{\small
\begin{tabularx}{\linewidth}{@{}lrrYYYYr@{}}
\toprule
\textbf{Run} & \textbf{D} & \textbf{Seeds} & \textbf{Top-1 (\%)} & \textbf{Top-3 (\%)} & \textbf{Macro-F1 (\%)} & \textbf{CE loss} & \textbf{Phút} \\
\midrule
\SummaryRows
\bottomrule
\end{tabularx}}
\textit{Cách đọc:} mean $\pm$ population std của 3 seeds. Std của accuracy/F1 được biểu diễn bằng điểm phần trăm. CE loss là evaluation loss không có label smoothing. Thời gian là tổng training loop của 3 seeds, không gồm preprocessing, split, xuất báo cáo hoặc nén ZIP.

\section*{Bảng 2. Các phép so sánh có kiểm soát}
\begin{tabularx}{\linewidth}{@{}L{43mm}Yr@{}}
\toprule
\textbf{So sánh} & \textbf{Thành phần được đo} & \textbf{$\Delta$ Macro-F1 (điểm)} \\
\midrule
\ComparisonsRows
\bottomrule
\end{tabularx}

\section*{Nhận xét ở thời điểm báo cáo}
F0-control là mốc đối chứng cho F1 vì cùng tắt wrist scale/translation. So F1 trực tiếp với F0 sẽ thay đồng thời hệ tọa độ và augmentation, khiến khó quy kết nguồn cải thiện. Với kết quả hiện có, bỏ augmentation này cải thiện mean Macro-F1 trên cùng cohort; hiệu quả chưa đồng nhất ở từng seed và từng signer.

\textbf{PENDING} nghĩa là chưa có bộ artifacts đủ điều kiện đưa vào bảng, không phải điểm 0. Báo cáo chưa suy ra kết quả của các cấu hình chưa chạy/chưa nhận và chưa đưa test metrics vào lựa chọn feature.

\newpage
\pageheading{THIẾT KẾ THỰC NGHIỆM / 02}{Protocol và dữ liệu cố định}
\section*{Bảng 3. Những yếu tố giữ nguyên}
\begin{tabularx}{\linewidth}{@{}L{40mm}Y@{}}
\toprule
\textbf{Thành phần} & \textbf{Protocol} \\
\midrule
Dataset / vocabulary & Kaggle \code{asl-signs}; 20 glosses có thứ tự cố định; dùng chung selection, label map và frozen signer split. \\
Preprocessing & Trim đầu/cuối không có tay; giữ buffer 2 frames; ít nhất 3 valid hand frames; nội suy gap tối đa 3; resample $T=32$; dominant hand \code{auto}; epsilon $10^{-4}$. \\
Features nền & \code{features.set=hands\_pose}; hệ tọa độ, palm scale, tips và bends thay theo Bảng 5. \\
Model & GRU 2 layers, hidden size 128, unidirectional; \code{mean\_max} pooling; dropout 0,3; đầu ra 20 logits. \\
Optimizer / regularization & Adam; batch 32; learning rate 0,001; weight decay 0,0001; label smoothing 0,1; gradient clipping 1,0. \\
Scheduler / stopping & ReduceLROnPlateau, factor 0,5, patience 5; tối đa 60 epochs; early stopping patience 15, min delta 0,0001, tối thiểu 5 epochs. \\
Seeds / initialization & 42, 43, 44; mỗi model khởi tạo mới; \code{training.resume\_checkpoint=null}. \\
Runtime đo hiện có & Kaggle, Tesla T4; PyTorch 2.11.0+cu130; preload \code{device}; AMP tắt. Không dùng preset T4 khác protocol. \\
Checkpoint / feature selection & Best checkpoint theo lowest validation loss có smoothing; chọn feature theo mean Validation Macro-F1; tie-break Top-1 rồi CPU p95 nếu đã đo. \\
Test policy & \code{evaluation.evaluate\_test=false}; final test sau khi chốt cấu hình và checkpoint. \\
\bottomrule
\end{tabularx}

\section*{Bảng 4. Cohort của F0 và F0-control}
Cả hai xử lý 8.085 sequences: giữ 7.854 (97,14\%), loại 231 (2,86\%). Trong số loại, 122 chỉ có 2 valid hand frames và 109 chỉ có 1; không đủ ngưỡng 3. Các tập có đủ 20 glosses.

\begin{tabularx}{\linewidth}{@{}lrrY@{}}
\toprule
\textbf{Tập} & \textbf{Sequences} & \textbf{Signers} & \textbf{Vai trò} \\
\midrule
Train & 5.558 & 15 & Tối ưu trọng số, augmentation chỉ trong train. \\
Validation & 1.155 & 3 & Chọn checkpoint và so sánh feature. Signers: 28656, 32319, 36257. \\
Test giữ lại & 1.141 & 3 & Signers: 27610, 37779, 61333; chỉ ghi membership/counts, chưa evaluate trong campaign. \\
\bottomrule
\end{tabularx}

Signer sets không giao nhau. \CohortWarning{} F1--F4 phải rebind chính frozen split; nếu exclusions làm cohort đổi thì bổ sung matched-cohort evaluation trên giao các mẫu validation.

\section*{Thước đo và đơn vị}
Top-1 đo dự đoán đúng nhất; Top-3 đo nhãn đúng nằm trong ba dự đoán đầu. Macro-F1 là trung bình F1 của 20 lớp, để mỗi gloss có trọng số ngang nhau:
\[\operatorname{MacroF1}=\frac1{20}\sum_{c=1}^{20}\frac{2TP_c}{2TP_c+FP_c+FN_c},\qquad
\bar m=\frac13\sum_s m_s,\quad \sigma=\sqrt{\frac13\sum_s(m_s-\bar m)^2}.\]
Std mô tả biến thiên giữa seeds trên \emph{một} signer split; không phải confidence interval hay biến thiên qua nhiều split.

\newpage
\pageheading{BẢNG ABLATION CHÍNH / 03}{Cấu hình, lý do chạy và đầu ra}
\section*{Bảng 5. Ma trận cấu hình}
{\small
\begin{tabularx}{\linewidth}{@{}lYYYYcr@{}}
\toprule
\textbf{Run} & \textbf{Wrist reference} & \textbf{Palm scale} & \textbf{Tips} & \textbf{Bends} & \textbf{Scale/translate} & \textbf{D} \\
\midrule
F0 & image & frame & false & false & Bật & 160 \\
F0-control & image & frame & false & false & Tắt & 160 \\
F1 & torso\_xy & frame & false & false & Tắt & 162 \\
F2 & torso\_xy & clip\_median & false & false & Tắt & 162 \\
F3 & torso\_xy & clip\_median & true & false & Tắt & 192 \\
F4 & torso\_xy & clip\_median & true & true & Tắt & 212 \\
\bottomrule
\end{tabularx}}
{\small\raggedright Tên cột dùng các canonical keys: \code{features.wrist\_reference}, \code{features.palm\_scale},
\code{features.fingertip\_displacement}, \code{features.bend\_cosines} và \code{augmentation.scale\_translate.enabled}.\par}

\section*{Bảng 6. Tại sao cần từng run?}
{\small
\begin{tabularx}{\linewidth}{@{}L{23mm}YYL{32mm}@{}}
\toprule
\textbf{Run / người} & \textbf{Lý do và giả thuyết cần kiểm chứng} & \textbf{So sánh và đầu ra cần đọc} & \textbf{Kết quả hiện có} \\
\midrule
\textbf{F0 / A} & Tạo baseline đầy đủ hands\_pose với wrist trong hệ image và palm scale theo frame. Giữ augmentation gốc để định lượng chất lượng cấu hình nền trước thay đổi. & Tensor $(32,160)$; 3 model; validation metrics, confusion và per-signer. Là mốc cho F0-control. & Đã đủ 3 seeds; Macro-F1 \BaselineFone\%. \\
\midrule
\textbf{F0-control / A} & Tách tác dụng của wrist-only scale/translation khỏi tác dụng body-relative geometry. Phép biến đổi chỉ tác động wrist, không biến đổi đồng bộ torso; cần đối chứng trước F1. & Chỉ tắt scale/translate so với F0. D vẫn 160; đọc $\Delta$ Macro-F1 và consistency theo seed/gloss/signer. & Đã đủ 3 seeds; Macro-F1 \ControlFone\%; $\Delta$ \FoneDelta{} điểm. \\
\midrule
\textbf{F1 / B} & Thử wrist tương đối với tâm hai vai, chuẩn hóa theo độ rộng vai của clip. Giả thuyết: giảm phụ thuộc vị trí người trong image, giữ vị trí tay so với cơ thể. & So với F0-control. Thêm 2 geometry-support channels; tensor $(32,162)$. Đọc gains và tỷ lệ geometry hợp lệ. & \RunOneOutcome \\
\midrule
\textbf{F2 / B} & Thử palm scale bằng median của valid observations trong clip thay cho từng frame. Giả thuyết: giảm dao động handshape do nhiễu ước lượng scale. & So với F1; không thêm channels, D=162. Kiểm tra accuracy/F1 và các lỗi do bàn tay thay đổi hướng. & \RunTwoOutcome \\
\midrule
\textbf{F3 / C} & Bổ sung chuyển động cục bộ của 5 đầu ngón mỗi tay. Giả thuyết: phân biệt động tác ngón tay mà quỹ đạo wrist hoặc handshape tĩnh chưa mô tả rõ. & So với F2; thêm $2\times5\times3=30$ channels, D=192. Đọc per-gloss/confusion và chi phí tăng. & \RunThreeOutcome \\
\midrule
\textbf{F4 / C} & Bổ sung cosine giữa các đoạn xương ngón tay để biểu diễn độ gập. Giả thuyết: làm rõ handshape; tính bất biến với scale/rotation hữu ích khi landmarks có support. & So với F3; thêm $2\times5\times2=20$ channels, D=212. Đọc gains, degeneracy và chi phí. & \RunFourOutcome \\
\bottomrule
\end{tabularx}}
\textbf{Điều kiện chung:} flip (p=0,5; swap\_hands=false), jitter (p=0,5; sigma=0,01), rotate (p=0,5; tối đa $8^\circ$), shape\_scale (p=0,5; r=0,10) giữ nguyên; time\_stretch tắt. F0 scale/translate: p=0,7, scale $[0,85;1,15]$, translation tối đa 0,10 trên x/y.

\newpage
\pageheading{GIẢI THÍCH FEATURE / 04}{Cơ chế và ý nghĩa đầu ra}
\section*{Feature nền: vì sao $D=160$?}
\begin{tabularx}{\linewidth}{@{}YrY@{}}
\toprule
\textbf{Nhóm} & \textbf{D} & \textbf{Thông tin cung cấp} \\
\midrule
Local hand shapes & 126 & 2 tay $\times$ 21 landmarks $\times$ 3 tọa độ; trừ wrist và chia palm scale. \\
Wrists / relative wrist / velocity & 15 & 6 tọa độ wrist, 3 chênh lệch hai wrist, 6 sai phân theo thời gian. \\
Hand masks & 2 & Support của từng hand slot; phân biệt giá trị thiếu với tọa độ thực. \\
Upper-body pose & 17 & 12 tọa độ hai vai/hai khuỷu tay chuẩn hóa, 3 torso velocity, 1 shoulder scale, 1 pose mask. \\
\bottomrule
\end{tabularx}
Tổng $126+15+2+17=160$. Output feature mỗi clip là float32 $(T,D)$, model nhận $(B,32,D)$ và trả $(B,20)$ logits. Top-1/Top-3 lấy từ 20 class probabilities.

\section*{F1: wrist trong hệ torso\_xy}
Gọi $\mathbf c_{t,xy}$ là tâm hai vai, $\ell$ là median khoảng cách hai vai trên các valid observed frames, $\mathbf w_{t,h}$ là wrist của tay $h$. Tọa độ được tính:
\[\widetilde{\mathbf w}_{t,h,xy} =
\frac{\mathbf w_{t,h,xy}-\mathbf c_{t,xy}}{\ell},\qquad
\widetilde w_{t,h,z}=0.\]
Wrist/relative wrist/velocity z bằng 0 theo contract torso\_xy. Hai geometry-support channels cho biết khi hand, pose và scale đủ điều kiện sử dụng. Unsupported geometry và transition không hợp lệ được zero-fill; không suy ra chuyển động qua gap thiếu support. Vì vậy F1 đo cả đổi representation \emph{và} thêm support channels, chưa tách đóng góp riêng của mỗi phần.

\section*{F2: palm scale ổn định theo clip}
\[s_{t,h}=\|\mathbf p_{t,h,9}-\mathbf p_{t,h,0}\|_2,\quad
s_h^{\mathrm{clip}}=\operatorname{median}_{t\in\mathcal V_h}s_{t,h},\quad
\widetilde{\mathbf p}_{t,h,j}=\frac{\mathbf p_{t,h,j}-\mathbf p_{t,h,0}}{s_h^{\mathrm{clip}}}.\]
Median lấy từ observed valid, non-degenerate frames trước nội suy/resampling theo feature pipeline. F1 dùng $s_{t,h}$ còn F2 dùng $s_h^{\mathrm{clip}}$; số channels không đổi. Cần kiểm chứng trade-off giữa giảm nhiễu scale và giữ biến thiên do foreshortening/hướng bàn tay; chưa thể khẳng định F2 tốt hơn F1.

\section*{F3: fingertip displacement}
Với các landmark indices $j\in\{4,8,12,16,20\}$, lấy sai phân local normalized fingertip giữa hai bước liên tiếp:
\[\Delta\widetilde{\mathbf p}_{t,h,j}=\widetilde{\mathbf p}_{t,h,j}-\widetilde{\mathbf p}_{t-1,h,j}.\]
Bước đầu bằng 0; chỉ giữ khi cả hai endpoints và transition hợp lệ. Đây là displacement trên chuỗi đã resample, không phải velocity theo giây. F3 cung cấp 30 channels mô tả articulation motion trực tiếp, nhưng có thể nhạy với nhiễu landmarks.

\section*{F4: bend cosines}
Hai cặp đoạn xương liên tiếp cho mỗi finger chain tạo cosine:
\[b=\operatorname{clip}\!\left(\frac{\mathbf u\cdot\mathbf v}{\|\mathbf u\|\|\mathbf v\|},-1,1\right).\]
20 channels: 2 tay $\times$ 5 ngón $\times$ 2 cặp. Theo vector convention, ngón thẳng gần $+1$, góc vuông gần 0; xương degenerate/unsupported cho 0. Cosine bổ sung inductive bias về articulation, không phải thông tin độc lập hoàn toàn với local landmarks.

\newpage
\pageheading{CHI TIẾT TRAINING / 05}{Kết quả từng seed và phép đối chứng}
\section*{Bảng 7. Checkpoint và validation từng model}
{\footnotesize\renewcommand{\arraystretch}{1.19}
\begin{tabularx}{\linewidth}{@{}lrrrYYYYr@{}}
\toprule
\textbf{Run} & \textbf{Seed} & \textbf{Best ep.} & \textbf{Tổng ep.} & \textbf{Top-1 \%} & \textbf{Top-3 \%} & \textbf{F1 \%} & \textbf{CE loss} & \textbf{Giây} \\
\midrule
\SeedsRows
\bottomrule
\end{tabularx}}
Best epoch là vị trí 1-based có validation loss nhỏ nhất trong history. Tổng epochs là độ dài history; thời gian lấy từ \code{training\_time\_sec}. Cấu hình chưa có kết quả giữ từng seed ở PENDING.

\section*{Bảng 8. Chênh lệch F0-control trừ F0 theo seed}
\begin{tabularx}{\linewidth}{@{}rYYYY@{}}
\toprule
\textbf{Seed} & \textbf{$\Delta$ Top-1} & \textbf{$\Delta$ Top-3} & \textbf{$\Delta$ Macro-F1} & \textbf{$\Delta$ CE loss} \\
\midrule
\PairedRows
\bottomrule
\end{tabularx}
Accuracy/F1 chênh lệch theo điểm phần trăm; CE loss giữ đơn vị gốc. Macro-F1 cải thiện ở \PositiveSeeds/3 seeds. Mean cải thiện không chứng minh mọi lần khởi tạo đều tốt hơn; chưa làm significance test và chưa có nhiều signer splits.

\section*{Hai loại loss cần đọc riêng}
\code{history.json} dùng CrossEntropyLoss với label smoothing 0,1: best validation loss của F0 theo seeds là 1,2783 / 1,3002 / 1,2792; của F0-control là 1,2842 / 1,2711 / 1,2504. Đây là loss chọn checkpoint, không phải CE loss trong Bảng 1 và Bảng 7.

\code{val\_report.json} dùng CrossEntropyLoss không smoothing, sau khi nạp best weights. Việc report loss khoảng 0,85--0,90 trong khi history khoảng 1,25--1,30 phù hợp với hai objective khác nhau; không thay một loại bằng loại còn lại.

\textbf{Giới hạn khi đọc thời gian:} F0 và F0-control chạy trên cùng loại Tesla T4, nhưng số epochs thực tế khác nhau. Tổng training time thấp hơn không đủ chứng minh throughput mỗi epoch tốt hơn. Chưa đo end-to-end CPU latency p50/p95, peak memory hoặc realtime FPS.

\newpage
\pageheading{CHẨN ĐOÁN VALIDATION / 06}{Theo signer và từng gloss}
\section*{Bảng 9. Top-1 theo validation signer (\%)}
\begin{tabularx}{\linewidth}{@{}lrYYr@{}}
\toprule
\textbf{Signer} & \textbf{Mẫu} & \textbf{F0: mean $\pm$ std} & \textbf{F0-control: mean $\pm$ std} & \textbf{$\Delta$ (điểm)} \\
\midrule
\SignersRows
\bottomrule
\end{tabularx}
Mean/std tính từ accuracy từng signer trong ba \code{val\_report.json}; không tính bằng cách gộp signers với trọng số ngang nhau. Macro-F1 toàn cohort tăng không có nghĩa mọi signer đều tăng.

\section*{Bảng 10. F1 từng gloss trên cùng 1.155 mẫu}
{\small\renewcommand{\arraystretch}{1.12}
\begin{tabularx}{\linewidth}{@{}lrYYr@{}}
\toprule
\textbf{Gloss} & \textbf{Val mẫu} & \textbf{F0 F1 (\%)} & \textbf{F0-control F1 (\%)} & \textbf{$\Delta$ (điểm)} \\
\midrule
\GlossesRows
\bottomrule
\end{tabularx}}
F1 của từng gloss được tính từ confusion matrix của \emph{mỗi seed}, sau đó lấy mean qua ba seeds; không gộp confusion matrices trước khi tính F1. Số mẫu là support một validation cohort, không nhân ba.

\textbf{Nhóm tăng nhiều nhất:} \GlossGains. \textbf{Nhóm giảm nhiều nhất:} \GlossDeclines. Những biến động này là quan sát thực nghiệm, chưa xác lập cơ chế nguyên nhân.

\textbf{Điểm cần theo dõi khi có F1--F4:} awake/wake và mouse/pretend có F1 thấp hơn nhiều gloss khác. Kiểm tra confusion cụ thể và tỷ lệ geometry/transition support để xem feature mới có xử lý các lỗi đó hay chỉ đổi phân bố lỗi. Hiện chưa có chẩn đoán của F1--F4.

\newpage
\pageheading{ARTIFACTS VÀ TÁI LẬP / 07}{Đầu ra, cohort và danh tính thí nghiệm}
\section*{Bảng 11. Đầu ra kiểm tra theo cấu hình}
{\small
\begin{tabularx}{\linewidth}{@{}lYYYrrl@{}}
\toprule
\textbf{Run} & \textbf{Train mẫu/người} & \textbf{Val mẫu/người} & \textbf{Test mẫu/người} & \textbf{Loại} & \textbf{Lớp} & \textbf{Artifacts} \\
\midrule
\CohortRows
\bottomrule
\end{tabularx}}
Test ở bảng này chỉ là frozen split metadata. Các kết quả model được xác nhận thuộc \code{split\_identity=val}; không có \code{test\_report.json} trong seed artifacts đã kiểm tra.

\section*{Mỗi run cần nộp những đầu ra nào?}
\begin{enumerate}
\item \textbf{Feature và data audit:} preprocessing schema/identity; manifest; exclusions theo reason/gloss/signer; split CSVs và label map; sample membership. Runtime cache chứa tensor $(32,D)$ và validity metadata; ZIP giữ audit, không đóng gói toàn bộ cache lớn.
\item \textbf{Ba model và metrics:} \code{seed\_42}, \code{seed\_43}, \code{seed\_44}; mỗi seed có \code{best\_model.pt}, \code{history.json}, \code{train\_metrics.csv}, \code{val\_report.json}. Report gồm Top-1/Top-3/Macro-F1/loss, confusion matrix, per-participant accuracy.
\item \textbf{Provenance và logs:} \code{metadata.json}, \code{effective\_config.yaml}, input source/config/split/selection snapshots, stage commands và exit codes; \code{preflight.log}, \code{preprocess.log}, \code{split.log}, \code{train.log}.
\item \textbf{Handoff:} \code{multi\_seed\_summary.json}, attempt ZIP, \code{RESULTS.md}/\code{RESULTS.json}. Báo cáo LaTeX chỉ đưa số đã xác minh; thiếu artifacts hoặc chưa đủ ba seeds không tính mean/std chính thức.
\end{enumerate}

\section*{Bảng 12. Danh tính của F0/F0-control đã xác minh}
{\footnotesize
\begin{tabularx}{\linewidth}{@{}L{40mm}Y@{}}
\toprule
\textbf{Trường} & \textbf{Giá trị} \\
\midrule
Source commit chung & \SourceCommit \\
Checkout commit lúc train & \CheckoutCommit \\
Config snapshot SHA256 & \ConfigHash \\
Frozen split SHA256 & \SplitHash \\
Ordered selection SHA256 & \SelectionHash \\
Preprocess identity & \PreprocessHash \\
\bottomrule
\end{tabularx}}
Source SHA và checkout SHA có thể khác vì report commits; runner xác minh source/configs/scripts/requirements khớp freeze. F0 và F0-control có cùng preprocess identity vì augmentation thay trong train, không đổi cached feature representation.

\textbf{Attempts thực tế:} F0: \code{\BaselineAttempt}; F0-control: \code{\ControlAttempt}. Giờ chạy quy đổi +07:00: F0 \BaselineWindow; F0-control \ControlWindow. Bản gốc ghi UTC trong metadata; không đổi timestamp trong tên file.

\newpage
\pageheading{GIỚI HẠN VÀ CẬP NHẬT / 08}{Kết luận tạm thời và bước hoàn thiện}
\section*{Những kết luận có thể dùng hiện tại}
F0-control cải thiện mean Validation Macro-F1 \FoneDelta{} điểm trên cùng cohort so với F0, và là baseline hợp lệ cho phép so sánh body-relative wrist của F1. Kết luận này giới hạn ở 20 glosses, seeds 42/43/44 và frozen signer split hiện tại. Chưa chọn winner toàn campaign; F4 không mặc định thắng vì nhiều channels hơn.

\section*{Giới hạn và phần cần bổ sung}
\begin{itemize}
\item \textbf{Selection provenance:} vocabulary kế thừa từ \code{islr\_phase2\_artifacts\_15.10.2026}. Shared preparation ghi rõ historical selection/test usage chưa được chứng nhận untouched. Không mô tả study như một đánh giá final test hoàn toàn độc lập chưa từng chịu ảnh hưởng của campaign trước.
\item \textbf{Sequential ablation:} các feature cộng dồn. F2--F4 đo lợi ích có điều kiện trên cấu hình trước; chưa tách mọi interaction, chưa factorial ablation. F1 cũng thêm support channels cùng geometry.
\item \textbf{Model capacity:} tăng D làm tăng tham số input của GRU dù hidden size/layers giữ nguyên. Theo kiến trúc hiện tại, số tham số tính lý thuyết là 215.572 (D160), 216.340 (D162), 227.860 (D192), 235.540 (D212). Cần ghi nhận capacity cùng representation; chưa có capacity-matched control.
\item \textbf{Tổng quát hóa:} ba seeds trên một split chưa đánh giá biến thiên giữa các tập signers. Chưa có confidence intervals, significance test, robustness ngoài dataset hoặc đo CPU p95/peak memory.
\item \textbf{Chờ bổ sung:} \PendingSummary{} Geometry/palm/transition/bend diagnostics cần review riêng; architecture/T/label-smoothing ablations và final test chưa đo. Không điền giá trị kỳ vọng vào bảng.
\end{itemize}

\section*{Cập nhật báo cáo sau các run còn lại}
\begin{enumerate}
\item B/C chạy seeds 42/43/44 từ initialization mới, dùng đúng shared source/config/selection/split và tắt test evaluation; nộp attempt ZIP cùng logs/audits.
\item Giải nén vào \path{artifacts/ablation-study/<run>/<attempt>/}. Nếu có nhiều attempts cho một run, chỉ rõ các attempts đã chọn; không tự chọn latest hoặc best seed.
\item Từ repo root chạy \path{reports/ablation-study/latex/build_report.ps1}. Script gọi công cụ tổng hợp hiện có, kiểm tra artifacts, cập nhật bảng LaTeX và compile hai lần thành \code{ablation\_study.pdf}. Không chạy lại training.
\item Review cohort/exclusions, per-gloss/per-signer và các delta trước khi chốt winner. Sau đó mới chọn checkpoint và chạy final test theo protocol riêng; cập nhật kết luận bằng số thực tế.
\end{enumerate}

\section*{Nguồn nội bộ dùng cho báo cáo}
{\small
\begin{enumerate}
\item \path{reports/ablation-study/RESULTS.json}: trạng thái, aggregate/per-seed metrics, provenance của các runs đã nhận.
\item \path{reports/ablation-study/README.md} và \path{reports/ablation-study/shared/config.yaml}: ma trận sáu cấu hình, phân công và common protocol.
\item \path{artifacts/ablation-study/shared/preparation_provenance.json}: nguồn vocabulary và giới hạn historical test-use.
\item Attempt của F0/F0-control: \code{metadata.json}, \code{cache\_audit/}, \code{split\_snapshot/}, \code{history.json}, \code{val\_report.json}; xác minh bằng \path{reports/ablation-study/tools/summarize_results.py}.
\item \path{src/features/schema.py}, \path{src/features/normalization.py}, \path{src/features/pipeline.py}: layout, normalization, support và derived features.
\item \path{src/data/augmentation.py}, \path{src/data/augmentation_gpu.py}: wrist scale/translation và augmentation contract.
\item \path{src/models/gru.py}, \path{src/training/train.py}, \path{src/training/evaluate.py}: model, checkpoint selection và hai objective loss.
\end{enumerate}}
\textit{Lưu ý nguồn:} các mẫu \code{runs/*.md} và \code{REPORT.md} cũ vẫn có mục PENDING. Số trong tài liệu này lấy từ numerical results và artifacts thực tế; không dùng trạng thái cũ để ghi đè kết quả đã hoàn tất.
\end{document}
